\documentclass[11pt,a4paper]{article}

\usepackage[utf8]{inputenc}
\usepackage[T1]{fontenc}
\usepackage{amsmath,amssymb,amsfonts}
\usepackage{physics}
\usepackage{hyperref}
\usepackage{booktabs}
\usepackage{geometry}
\usepackage{enumitem}
\usepackage{xcolor}
\usepackage{graphicx}
\usepackage{subcaption}
\usepackage{multirow}

\geometry{margin=1in}
\hypersetup{colorlinks=true, linkcolor=blue, urlcolor=blue, citecolor=blue}

\title{Comparison of VQE, VQITE, and TTITE\\
  for Charmonium Spectroscopy\\[0.5em]
  \large Apples-to-Apples Benchmark on the Woloshyn Hamiltonian}

\author{quarkonium-simulations project}
\date{\today}

\begin{document}
\maketitle

\begin{abstract}
We compare three quantum-computing methods---the Variational Quantum
Eigensolver (VQE), Variational Quantum Imaginary-Time Evolution (VQITE),
and Trotter--Taylor Imaginary-Time Evolution (TTITE)---on the same
2-qubit charmonium Hamiltonian with spin-dependent interactions.
All quantum operations run through Qiskit's \texttt{AerSimulator}:
energies are measured via shot-based Pauli circuits, VQITE gradients
and metric tensors via parameter-shift and overlap circuits, and TTITE
evolution via actual LCU circuits with ancilla postselection.
Under depolarizing noise, VQE proves most resilient (shallow circuits),
VQITE degrades due to noisy gradient estimation, and TTITE collapses
to the maximally mixed state ($\sim\!2250$ sequential noisy gates).
\end{abstract}

\tableofcontents
\newpage

%=============================================================================
\section{Setup}
%=============================================================================

\subsection{The Hamiltonian}

We use the Woloshyn charmonium Hamiltonian (arXiv:2301.10828v2), which
describes a charm--anticharm system with the Cornell potential plus a
spin-dependent contact term:
\begin{equation}
  V(r) = -\frac{a}{r} + b\,r + V_s(r)\,\vec{S}_c \cdot \vec{S}_{\bar{c}}
\end{equation}
The Hamiltonian is expanded in a harmonic oscillator basis truncated
to $N=4$ radial states, yielding a $4 \times 4$ matrix encoded on
\textbf{2 qubits} via direct state mapping:
$\ket{00} \leftrightarrow n\!=\!0$, \;
$\ket{01} \leftrightarrow n\!=\!1$, \;
$\ket{10} \leftrightarrow n\!=\!2$, \;
$\ket{11} \leftrightarrow n\!=\!3$.

Three spin channels are studied:
${}^1S_0$ ($\eta_c$, spin singlet),
${}^3S_1$ ($J/\psi$, spin triplet), and
${}^1P_1$ ($h_c$, $P$-wave singlet).
The Pauli decomposition of each channel has 10 terms:
\begin{equation}
  H = c_0\,II + c_1\,IZ + c_2\,ZI + c_3\,ZZ + c_4\,IX + c_5\,XI
    + c_6\,ZX + c_7\,XZ + c_8\,XX + c_9\,YY
\end{equation}

\subsection{The Three Methods}

\paragraph{VQE (Variational Quantum Eigensolver).}
Minimizes the energy expectation value over a parameterized ansatz:
\begin{equation}
  E(\vec{\theta}) = \mel{\psi(\vec{\theta})}{H}{\psi(\vec{\theta})} \geq E_0
\end{equation}
using the COBYLA classical optimizer. Each energy evaluation measures
all 10 Pauli terms via shot-based circuits on \texttt{AerSimulator}
(8\,192 shots per term).

\paragraph{VQITE (Variational Quantum Imaginary-Time Evolution).}
Evolves variational parameters along the imaginary-time gradient:
\begin{equation}
  \sum_j A_{ij}\,\dot{\theta}_j = C_i, \qquad
  C_i = -\frac{\partial E}{\partial\theta_i}, \quad
  A_{ij} = \text{Re}\!\left[\frac{\partial\bra{\psi}}{\partial\theta_i}
    \frac{\partial\ket{\psi}}{\partial\theta_j}\right]
\end{equation}
The gradient $C_i$ is computed via the parameter-shift rule on shot-based
energy circuits. The quantum metric tensor $A_{ij}$ is computed via
overlap circuits ($U(\theta_0)\,U(\theta')^\dagger$, measuring $P(\ket{00})$)
with finite differences. All quantum operations run on \texttt{AerSimulator}.

\paragraph{TTITE (Trotter--Taylor Imaginary-Time Evolution).}
Implements imaginary-time evolution directly (no ansatz) via:
\begin{equation}
  e^{-H\tau} \approx \left[\prod_{i=1}^m \bigl(\alpha_i\,I + \beta_i\,h_i\bigr)\right]^L
\end{equation}
where $\alpha_i$, $\beta_i$ are Taylor expansion coefficients and each
factor is applied via an ancilla-qubit LCU circuit. In the noiseless
case, each step runs on \texttt{AerSimulator(method='statevector')}
with proper postselection (ancilla$\,=\ket{0}$). In the noisy case,
all $\sim\!2250$ LCU steps are chained into a single continuous circuit
with ancilla \texttt{reset} between steps, run on
\texttt{AerSimulator(method='density\_matrix', noise\_model=...)}.
The final energy is computed via $\text{Tr}(\rho H)$ from the
output density matrix.

\subsection{Shared Ansatz}

Both VQE and VQITE use the same hardware-efficient ansatz:
\begin{equation}
  \begin{aligned}
  &q_0: \quad R_Y(\theta_0) - \bullet - \\
  &q_1: \quad R_Y(\theta_1) - \oplus - R_Y(\theta_2)
  \end{aligned}
\end{equation}
With 3 parameters, this circuit spans all real superpositions of the 4 basis
states (3 free amplitudes after normalization). Circuit depth: 3 layers,
1 CNOT gate.

\subsection{Circuit Execution Summary}

Every quantum computation runs on \texttt{AerSimulator}---no numpy matrix
operations are used for energy calculation:

\begin{center}
\small
\begin{tabular}{@{}l l l@{}}
\toprule
Method & Operation & Circuit type \\
\midrule
VQE & Energy evaluation & Shot-based Pauli measurement \\
VQE & Optimization & Classical (COBYLA) \\
\midrule
VQITE & Energy gradient $C_i$ & Parameter-shift on shot-based circuits \\
VQITE & Metric tensor $A_{ij}$ & Overlap circuits ($U U'^\dagger$) \\
\midrule
TTITE (noiseless) & Evolution step & LCU circuit on statevector backend \\
TTITE (noisy) & Full evolution & Single continuous circuit, density matrix \\
TTITE & Final energy & $\text{Tr}(\rho H)$ (noisy) / shot-based (noiseless) \\
\bottomrule
\end{tabular}
\end{center}


%=============================================================================
\section{Noiseless Results}
%=============================================================================

\subsection{Ground-State Energy}

All three methods are run with shot noise only (no gate errors).
VQE and VQITE use 8\,192 shots per Pauli term; TTITE uses statevector
simulation for the evolution and shot-based measurement for the final energy.

\begin{center}
\small
\begin{tabular}{@{}l l r r r r r@{}}
\toprule
Channel & Method & $E_0$ [fm$^{-1}$] & Exact & Error & Fidelity & Time [s] \\
\midrule
1S0 & VQE   & 0.4375 & 0.3913 & $+$0.046 & 0.987 & 14 \\
1S0 & VQITE & 0.4192 & 0.3913 & $+$0.028 & 0.998 & 144 \\
1S0 & TTITE & 0.3560 & 0.3913 & $-$0.035 & 1.000 & 0.8 \\
\midrule
3S1 & VQE   & 0.7879 & 0.7522 & $+$0.036 & 0.986 & 19 \\
3S1 & VQITE & 0.9687 & 0.7522 & $+$0.217 & 0.925 & 141 \\
3S1 & TTITE & 0.7452 & 0.7522 & $-$0.007 & 1.000 & 0.8 \\
\midrule
1P1 & VQE   & 2.8041 & 2.7813 & $+$0.023 & 0.984 & 21 \\
1P1 & VQITE & 3.1544 & 2.7813 & $+$0.373 & 0.929 & 141 \\
1P1 & TTITE & 2.7819 & 2.7813 & $+$0.001 & 1.000 & 0.8 \\
\bottomrule
\end{tabular}
\end{center}

\textbf{Key observations:}
\begin{itemize}[nosep]
  \item \textbf{TTITE achieves the highest fidelity} ($>0.9998$) with
    sub-percent energy errors. Its evolution is deterministic (no optimizer
    noise) and uses statevector simulation for proper postselection.
  \item \textbf{VQE} converges reliably on 1S0 and 1P1 (errors $\sim$0.02--0.05)
    but sometimes gets trapped by COBYLA in shot-noise landscapes (3S1 error 0.04).
  \item \textbf{VQITE} struggles on 3S1 and 1P1 (errors 0.2--0.4). The
    circuit-based metric tensor (overlap circuits with finite differences)
    introduces significant noise that degrades the natural gradient updates.
  \item \textbf{Wall-clock time}: TTITE is fastest ($<$1\,s), VQE is moderate
    (15--20\,s), VQITE is slowest ($\sim$140\,s) due to $\sim$30 circuit
    evaluations per step $\times$ 40 steps.
\end{itemize}

\begin{figure}[htbp]
\centering
\includegraphics[width=0.85\textwidth]{energy_summary.png}
\caption{Noiseless ground-state energy comparison across all channels and methods.}
\label{fig:energy_summary}
\end{figure}

\begin{figure}[htbp]
\centering
\includegraphics[width=0.85\textwidth]{convergence_comparison.png}
\caption{Convergence traces. VQE: energy vs optimizer iteration;
  VQITE: energy vs imaginary-time step; TTITE: energy vs $\tau$.
  Dashed lines mark the exact eigenvalue.}
\label{fig:convergence}
\end{figure}

\begin{figure}[htbp]
\centering
\includegraphics[width=0.85\textwidth]{circuit_cost.png}
\caption{Computation cost: wall-clock time (left) and number of
  iterations/steps (right).}
\label{fig:cost}
\end{figure}


%=============================================================================
\section{Excited States}
%=============================================================================

Excited states are computed via VQE with the penalty method:
\begin{equation}
  E_{\text{eff}}(\vec{\theta}) = \mel{\psi(\vec{\theta})}{H}{\psi(\vec{\theta})}
  + \alpha \sum_{j<k} \bigl|\braket{\phi_j|\psi(\vec{\theta})}\bigr|^2
\end{equation}
with $\alpha = 10$. The overlaps $|\!\braket{\phi_j|\psi}\!|^2$ are measured
via the swap test circuit (ancilla + controlled-SWAPs). VQITE excited
states are skipped due to excessive circuit overhead (swap test in every
gradient evaluation $\times$ every VQITE step). TTITE does not support
excited states.

\begin{center}
\small
\begin{tabular}{@{}l c r r r@{}}
\toprule
Channel & State & VQE Energy & Exact & Error \\
\midrule
1S0 & 0 & 0.398 & 0.391 & $+$0.007 \\
1S0 & 1 & 3.375 & 3.483 & $-$0.108 \\
1S0 & 2 & 5.980 & 5.643 & $+$0.336 \\
1S0 & 3 & 5.369 & 7.545 & $-$2.177 \\
\midrule
3S1 & 0 & 0.797 & 0.752 & $+$0.044 \\
3S1 & 1 & 3.454 & 3.633 & $-$0.179 \\
3S1 & 2 & 5.817 & 5.727 & $+$0.091 \\
3S1 & 3 & 5.947 & 7.646 & $-$1.699 \\
\midrule
1P1 & 0 & 2.790 & 2.781 & $+$0.008 \\
1P1 & 1 & 4.848 & 4.883 & $-$0.035 \\
1P1 & 2 & 7.901 & 6.756 & $+$1.145 \\
1P1 & 3 & 7.005 & 8.770 & $-$1.765 \\
\bottomrule
\end{tabular}
\end{center}

Ground and first excited states are reasonable (errors $<$0.2~fm$^{-1}$).
Higher excited states (states 2--3) show large errors: the swap test
provides imprecise overlap estimates, and accumulated penalty errors
from lower states propagate upward. This is a known limitation of the
penalty method under shot noise.

\begin{figure}[htbp]
\centering
\includegraphics[width=0.85\textwidth]{excited_states.png}
\caption{Energy levels from VQE (with penalty method) compared to exact
  diagonalization. Only VQE results are shown; VQITE excited states are
  too slow with circuit-based overlap measurements; TTITE does not
  support excited states.}
\label{fig:excited}
\end{figure}


%=============================================================================
\section{Noise Analysis}
%=============================================================================

\subsection{Noise Model}

All methods are run under depolarizing noise at rates $p \in \{0, 0.01, 0.05\}$
with 3 independent repeats per configuration. The noise model applies
depolarizing error to every gate: rate $p$ for single-qubit gates
($R_Y$, $X$, $H$, $S^\dagger$) and rate $p$ for two-qubit gates (CNOT).

\begin{itemize}[nosep]
  \item \textbf{VQE}: Every energy evaluation runs shot-based Pauli
    measurement circuits on a noisy \texttt{AerSimulator}. The COBYLA
    optimizer sees noisy cost function values but can partially adapt.
  \item \textbf{VQITE}: Energy gradient via parameter-shift rule on
    noisy circuits. Metric tensor via overlap circuits on the same noisy
    backend. Both gradient and metric tensor are corrupted by gate noise.
  \item \textbf{TTITE}: A single continuous circuit chains all
    $\sim\!2250$ LCU steps with ancilla \texttt{reset} between steps,
    run on the density-matrix backend with the noise model.
    The \texttt{reset} instruction does \emph{not} postselect---if the
    ancilla measures $\ket{1}$ (failed LCU), the corrupted work state
    carries forward, modeling realistic hardware behavior without
    repeat-until-success.
\end{itemize}

\subsection{Results}

\begin{center}
\small
\begin{tabular}{@{}l l r@{\;\;}c@{\;\;}l r@{}}
\toprule
Channel & Noise & \multicolumn{3}{c}{VQE} & Error \\
\midrule
\multirow{3}{*}{1S0}
  & 0\%  & $0.547$ & $\pm$ & $0.146$ & $+$0.16 \\
  & 1\%  & $0.636$ & $\pm$ & $0.183$ & $+$0.24 \\
  & 5\%  & $1.044$ & $\pm$ & $0.090$ & $+$0.65 \\
\midrule
\multirow{3}{*}{3S1}
  & 0\%  & $0.912$ & $\pm$ & $0.092$ & $+$0.16 \\
  & 1\%  & $1.016$ & $\pm$ & $0.098$ & $+$0.26 \\
  & 5\%  & $2.179$ & $\pm$ & $1.103$ & $+$1.43 \\
\midrule
\multirow{3}{*}{1P1}
  & 0\%  & $2.810$ & $\pm$ & $0.003$ & $+$0.03 \\
  & 1\%  & $2.902$ & $\pm$ & $0.011$ & $+$0.12 \\
  & 5\%  & $3.841$ & $\pm$ & $0.852$ & $+$1.06 \\
\bottomrule
\end{tabular}
\hspace{1em}
\begin{tabular}{@{}l l r@{\;\;}c@{\;\;}l r@{}}
\toprule
Channel & Noise & \multicolumn{3}{c}{VQITE} & Error \\
\midrule
\multirow{3}{*}{1S0}
  & 0\%  & $0.542$ & $\pm$ & $0.192$ & $+$0.15 \\
  & 1\%  & $2.157$ & $\pm$ & $1.679$ & $+$1.77 \\
  & 5\%  & $4.608$ & $\pm$ & $1.991$ & $+$4.22 \\
\midrule
\multirow{3}{*}{3S1}
  & 0\%  & $0.797$ & $\pm$ & $0.059$ & $+$0.04 \\
  & 1\%  & $1.147$ & $\pm$ & $0.080$ & $+$0.39 \\
  & 5\%  & $5.569$ & $\pm$ & $1.075$ & $+$4.82 \\
\midrule
\multirow{3}{*}{1P1}
  & 0\%  & $2.800$ & $\pm$ & $0.008$ & $+$0.02 \\
  & 1\%  & $2.914$ & $\pm$ & $0.007$ & $+$0.13 \\
  & 5\%  & $5.180$ & $\pm$ & $0.369$ & $+$2.40 \\
\bottomrule
\end{tabular}
\end{center}

\begin{center}
\small
\begin{tabular}{@{}l l r@{\;\;}c@{\;\;}l r l@{}}
\toprule
Channel & Noise & \multicolumn{3}{c}{TTITE} & Error & Note \\
\midrule
\multirow{3}{*}{1S0}
  & 0\%  & $0.382$ & $\pm$ & $0.003$ & $-$0.01 & \\
  & 1\%  & $4.266$ & $\pm$ & $0.000$ & $+$3.87 & $= \text{Tr}(H)/4$ \\
  & 5\%  & $4.266$ & $\pm$ & $0.000$ & $+$3.87 & $= \text{Tr}(H)/4$ \\
\midrule
\multirow{3}{*}{3S1}
  & 0\%  & $0.743$ & $\pm$ & $0.015$ & $-$0.01 & \\
  & 1\%  & $4.440$ & $\pm$ & $0.000$ & $+$3.69 & $= \text{Tr}(H)/4$ \\
  & 5\%  & $4.439$ & $\pm$ & $0.000$ & $+$3.69 & $= \text{Tr}(H)/4$ \\
\midrule
\multirow{3}{*}{1P1}
  & 0\%  & $2.787$ & $\pm$ & $0.003$ & $+$0.01 & \\
  & 1\%  & $5.798$ & $\pm$ & $0.000$ & $+$3.02 & $= \text{Tr}(H)/4$ \\
  & 5\%  & $5.798$ & $\pm$ & $0.000$ & $+$3.02 & $= \text{Tr}(H)/4$ \\
\bottomrule
\end{tabular}
\end{center}

\textbf{Key finding:} Under any nonzero gate noise, TTITE converges to
the energy of the maximally mixed state $\rho = I/d$:
\begin{equation}
  E_{\text{mixed}} = \text{Tr}\!\left(\frac{I}{4}\,H\right) = \frac{\text{Tr}(H)}{4}
\end{equation}
with \emph{zero variance} across repeats. This is because the $\sim\!2250$
sequential noisy gates drive the density matrix to $I/4$ regardless of
the initial state. The fidelity drops to $1/d = 0.25$ (random guessing).

\begin{figure}[htbp]
\centering
\includegraphics[width=0.95\textwidth]{noise_sweep.png}
\caption{Energy error vs depolarizing rate for all three methods, all
  channels. Error bars show $\pm 1\sigma$ over 3 repeats. TTITE (red)
  jumps immediately to $\text{Tr}(H)/4$ at any nonzero noise.}
\label{fig:noise}
\end{figure}

\begin{figure}[htbp]
\centering
\includegraphics[width=0.95\textwidth]{fidelity_vs_noise.png}
\caption{Ground-state fidelity vs depolarizing rate. VQE (blue) maintains
  high fidelity; VQITE (green) degrades with high variance; TTITE (red)
  drops to $1/4$ (maximally mixed) at any noise.}
\label{fig:fidelity}
\end{figure}


%=============================================================================
\section{Discussion}
%=============================================================================

\subsection{Accuracy vs Noise Trade-off}

The three methods occupy distinct points on the accuracy--noise trade-off
curve. In the noiseless limit, the ranking is:
\begin{equation*}
  \text{TTITE} \;(\sim\!0.001) \;\ll\;
  \text{VQE} \;(\sim\!0.03) \;<\;
  \text{VQITE} \;(\sim\!0.1)
\end{equation*}
where errors are in fm$^{-1}$. Under 1\% depolarizing noise:
\begin{equation*}
  \text{VQE} \;(\sim\!0.2) \;\ll\;
  \text{VQITE} \;(\sim\!1) \;\ll\;
  \text{TTITE} \;(\sim\!4)
\end{equation*}
The ranking completely inverts. TTITE's theoretical superiority (exact,
ansatz-free, guaranteed convergence) is irrelevant on noisy hardware
because its deep circuit ($\sim\!2250$ gates) accumulates errors
faster than the signal can converge.

\subsection{Why VQE Wins on Noisy Hardware}

VQE's resilience stems from three factors:
\begin{enumerate}[nosep]
  \item \textbf{Shallow circuits}: the ansatz has depth 3 (2 $R_Y$ layers
    + 1 CNOT). Each energy evaluation accumulates noise from only
    $\sim\!5$ gates, vs $\sim\!2250$ for TTITE.
  \item \textbf{Optimizer adaptability}: COBYLA adjusts its simplex to
    the noisy landscape, partially compensating for systematic bias.
  \item \textbf{Independent evaluations}: each energy measurement is a
    fresh circuit execution. Noise does not compound across iterations---only
    the \emph{current} evaluation is noisy.
\end{enumerate}

\subsection{Why VQITE Degrades More Than VQE}

VQITE uses the same shallow circuits as VQE for energy measurements, but
additionally requires the quantum metric tensor $A_{ij}$ at every step.
The circuit-based metric tensor (overlap circuits with finite differences)
is more sensitive to noise than energy measurements because:
\begin{itemize}[nosep]
  \item Small parameter shifts ($\varepsilon = 0.05$) amplify shot noise
    in the finite-difference formula.
  \item The metric tensor is a second-derivative quantity, inherently
    noisier than the first-derivative gradient.
  \item Errors in $A$ propagate through the linear solve
    $\dot{\theta} = A^{-1}C$, potentially inverting the update direction.
\end{itemize}

\subsection{Why TTITE Collapses to the Mixed State}

TTITE's continuous circuit has two independent failure modes:
\begin{enumerate}[nosep]
  \item \textbf{Gate noise accumulation}: with $\sim\!2250$ gate operations,
    even 0.1\% error per gate compounds to total state corruption.
    After $n$ depolarizing channels with rate $p$, the off-diagonal elements
    of $\rho$ decay as $(1-p)^n$. For $n=2250$ and $p=0.01$:
    $(1-0.01)^{2250} \approx e^{-22.5} \approx 10^{-10}$.
  \item \textbf{Failed postselection}: the ancilla \texttt{reset}
    does not postselect. When the ancilla measures $\ket{1}$
    (probability $\sim\!50\%$ per step), the ``wrong'' LCU branch
    ($\alpha I - \beta h_i$) contaminates the work state. Over many
    steps, both branches mix equally $\to$ maximally mixed state.
\end{enumerate}

\subsection{Practical Considerations}

\begin{center}
\begin{tabular}{@{}l c c c@{}}
\toprule
Criterion & VQE & VQITE & TTITE \\
\midrule
Ansatz required       & Yes & Yes & No \\
Ancilla qubits        & 0 & 0 & 1 \\
Circuit depth per eval & 3 & 3 & $\sim\!2250$ \\
Circuits per run      & $\sim\!500$ & $\sim\!1200$ & 1 (continuous) \\
Noise resilience      & \textbf{Best} & Moderate & Worst \\
Noiseless accuracy    & Moderate & Moderate & \textbf{Best} \\
Excited states        & Yes & Yes$^*$ & No \\
Convergence guarantee & No & Conditional & Yes (noiseless) \\
\bottomrule
\end{tabular}
\end{center}
{\small $^*$VQITE excited states are impractically slow with circuit-based
metric tensor + penalty overlaps.}

\subsection{Conclusion}

For near-term noisy quantum hardware, \textbf{VQE is the clear winner}:
its shallow circuits minimize noise exposure, and the classical optimizer
provides robustness to shot noise. VQITE offers stronger theoretical
convergence guarantees (monotonic energy decrease via natural gradient)
but is undermined by the circuit cost of computing the quantum metric
tensor---each VQITE step requires $\sim\!30$ circuit evaluations, and
the finite-difference metric tensor is fragile under noise.

TTITE is exact in principle and produces the best noiseless results, but
it is \textbf{fundamentally impractical on noisy hardware}. Its deep
continuous circuit and lack of postselection cause the quantum state to
decohere to the maximally mixed state under any nonzero noise. TTITE
is a \emph{fault-tolerant-era algorithm}: it will become competitive
only when quantum error correction can protect the $\sim\!2250$
sequential gate operations from decoherence.

\end{document}
